\chapter{Decisions, games, and dynamic programming}
\label{chap:decisiones}

\section{Choosing under uncertainty}
A decision combines controllable alternatives with states of the environment
that the analyst does not control. A payoff table organizes the payoff $a_{ij}$
from alternative $i$ in state $j$. When prior probabilities $p_j$ are known,
the expected-value criterion selects
\[
 i^*\in\arg\max_i\sum_j p_j a_{ij}.
\]
This criterion assumes comparable payoffs and risk neutrality. A solution
maximizing average profit may involve substantial losses in a particular state.
QSB does not automatically transform money into utility or incorporate the
decision maker's risk aversion.

\subsection{Payoff tables and information}
The file \path{dec_payoff.json} considers a launch with demand Alta (High,
probability $0.4$) or Baja (Low, $0.6$). The input's Spanish identifiers are
retained here so that they match the files and output: Lanzar means Launch,
Esperar means Wait, Positivo means Positive, and Negativo means Negative.
The payoffs are:
\begin{center}
\begin{tabular}{lrr}
\toprule
Decision & Alta & Baja\\
\midrule
Lanzar & 100 & $-40$\\
Esperar & 20 & 10\\
\bottomrule
\end{tabular}
\end{center}
Without a market study, Lanzar yields $0.4(100)+0.6(-40)=16$, compared with
14 for Esperar. With perfect information, Lanzar would be chosen in Alta and
Esperar in Baja: $0.4(100)+0.6(10)=46$. The expected value of perfect
information (EVPI) is $46-16=30$, an upper bound on the price of information
under this model.

The study produces a Positivo or Negativo indicator. The
\texttt{indicatorLikelihoods} matrix uses indicator rows and state columns:
$P(+\mid Alta)=0.8$ and $P(+\mid Baja)=0.2$. Each column represents a
distribution over possible indicators. These are not posterior probabilities.
The policy with the study is Lanzar if the result is Positivo and Esperar if
it is Negativo. Its expected value is $33.6$; hence EVSI (expected value of
sample information) is $33.6-16=17.6$. This amount does not deduct the cost
of the study. If the study cost 20, purchasing it would not be worthwhile
under the expected-value criterion.

\begin{lstlisting}
swift run qsb validate-decision-json manual/examples/dec_payoff.json
swift run qsb solve-decision-json manual/examples/dec_payoff.json --backend native
\end{lstlisting}
The output contains \texttt{priorExpectedValues}, \texttt{indicatorAnalyses},
\texttt{bestPriorDecision}, EVSI, and EVPI. Keys retain their English names.
The solver requires a probability greater than $10^{-12}$ for every indicator;
otherwise, solution fails rather than inventing a posterior.

\subsection{Bayesian updating}
Bayes' rule converts prior probability and likelihood into posterior probability:
\[
 P(s_j\mid e_k)=\frac{P(e_k\mid s_j)p_j}{\sum_\ell P(e_k\mid s_\ell)p_\ell}.
\]
The file \path{dec_bayes.json} separates this operation from any payoff table.
It uses the same states and study. QSB obtains $P(+)=0.44$, $P(-)=0.56$,
and the following posteriors:
\begin{center}
\begin{tabular}{lrr}
\toprule
Evidence & Alta & Baja\\
\midrule
Positivo & $0.727273$ & $0.272727$\\
Negativo & $0.142857$ & $0.857143$\\
\bottomrule
\end{tabular}
\end{center}
\begin{lstlisting}
swift run qsb validate-decision-json manual/examples/dec_bayes.json
swift run qsb solve-decision-json manual/examples/dec_bayes.json
\end{lstlisting}
\texttt{likelihoods} again uses outcome rows and state columns. Changing the
order of names requires changing the corresponding rows or columns as well.
Probabilities must be finite, nonnegative, and consistent with the distributions
represented. Bayes updates beliefs; by itself it does not say which action to take.

\subsection{Trees and backward induction}
A tree makes the temporal order of decisions and observations explicit.
\texttt{rootID} identifies the root, each node has a unique \texttt{id}, and
its \texttt{childIDs} identify successors. The types are \texttt{decision},
\texttt{chance}, and \texttt{terminal}. Chance-branch probabilities are stored
in the child nodes, rather than in a list on the parent.

Backward induction assigns a terminal its payoff. A decision node adds its
local payoff to the maximum value of its children; a chance node adds its
local payoff to the expected value of its children. Local payoffs accumulate
along the path: do not repeat a cost in the leaves if it was already entered
in an ancestor.

In \path{dec_tree.json}, Invertir (Invest) costs 5 and yields 40 with
probability $0.6$ or $-10$ with probability $0.4$. Its value is
$-5+0.6(40)+0.4(-10)=15$; Conservar (Retain) yields 12. The root value is
15, and the policy selects Invertir.
\begin{lstlisting}
swift run qsb validate-decision-json manual/examples/dec_tree.json
swift run qsb solve-decision-json manual/examples/dec_tree.json
\end{lstlisting}
\begin{figure}[H]
\centering
\begin{tikzpicture}
\node[rectangle,draw,minimum size=6mm,label=above:{\small Elegir}] (R) at (0,0) {};
\node[qsb node,label=above:{\small Invertir},label=below:{$-5$}] (I) at (4,1) {};
\node[qsb terminal] (K) at (4,-1.5) {Conservar\\12};
\node[qsb terminal] (E) at (8,2) {Exito\\40};
\node[qsb terminal] (F) at (8,-0.2) {Fracaso\\$-10$};
\draw[qsb selected] (R) -- (I);
\draw[qsb alternative] (R) -- (K);
\draw[qsb arc] (I) -- node[qsb label,above] {$0.6$} (E);
\draw[qsb arc] (I) -- node[qsb label,below] {$0.4$} (F);
\end{tikzpicture}

\caption{Decision tree: the square represents a choice and the circle represents chance. The thick branch is the selected policy. Leaves show payoffs; the local investment cost is added once. Names retain the JSON identifiers.}
\label{fig:decision-tree}
\end{figure}

Review \texttt{nodeValues} for each node and \texttt{policy} for the decisions.
The method maximizes payoffs: to minimize costs, you can use their negatives,
explaining the convention. QSB warns about and normalizes chance-child
probabilities whose positive sum differs from one, for compatibility with
rounded inputs. This repair does not replace a review of the data. Avoid cycles,
nonexistent identifiers, and inconsistent probabilities.

\section{Zero-sum games}
In a zero-sum game, the row player's payoff is the column player's loss. The
first maximizes and the second minimizes. Mixed strategies are probability
distributions over actions; they do not represent fractional quantities of an
action in a single play.

For a matrix $A$, QSB solves equivalent linear programs. The row player seeks
$\max v$ subject to $A^Tp\geq v\mathbf{1}$, $\sum_i p_i=1$, and $p_i\geq0$.
The column player solves the dual minimization problem. This assumes simultaneous
choices and known payoffs. \path{dec_game.json} uses
\[
 A=\begin{pmatrix}1&-1\\-1&1\end{pmatrix}.
\]
Both players assign probability $0.5$ to A and B, and the game's value is
zero. The numerical output may show $4.44\times10^{-16}$ because of rounding.
\begin{lstlisting}
swift run qsb validate-decision-json manual/examples/dec_game.json
swift run qsb solve-decision-json manual/examples/dec_game.json
\end{lstlisting}
Read \texttt{rowStrategy}, \texttt{columnStrategy}, and \texttt{value}. The
native method uses the shared linear solver. It does not interpret this solution
as an equilibrium of a general non-zero-sum game or a sequential game.

\section{Dynamic programming: state, stage, and action}
Dynamic programming divides a problem into stages. The state summarizes the
information needed to evaluate future decisions; an action changes that state
and produces a local cost or return. The principle of optimality allows an
optimal policy to be built from subproblems. A state omitting relevant
information breaks that principle.

QSB implements three discrete, deterministic models: bounded knapsack,
stage-by-stage acyclic routing, and production with inventory. Inputs use
\texttt{kind} and \texttt{model}. Outputs include a typed result,
\texttt{assumptions}, and \texttt{trace}. Each trace row describes a stage,
state, action, next state, and \emph{local} value; this value is not necessarily
the subproblem's optimal value function.

\subsection{Bounded knapsack}
Each item $i$ consumes $w_i$ units of capacity and contributes $r_i$ per unit;
at most $b_i$ units are available. The formulation is
\[
 \max\sum_i r_i x_i,\qquad \sum_i w_i x_i\leq C,
 \qquad x_i\in\{0,\ldots,b_i\}.
\]
A recurrence with remaining capacity $c$ is
$F_i(c)=\max_q\{q r_i+F_{i+1}(c-qw_i)\}$, over available quantities that fit.
Capacity and consumption are integers; returns may be finite real numbers.
Do not convert continuous physical measurements to integers without justifying
the discretization unit.

In \path{dec_knapsack.json}, $C=7$ and the items $(b_i,w_i,r_i)$ are
A $(2,3,5)$, B $(3,2,3)$, and C $(1,4,7)$. QSB selects one A and one C:
capacity 7 and return 12. The alternative of one A and two B also uses 7 but
returns 11. The selection encodes quantities; the trace also includes the
decision to take zero units of B.
\begin{lstlisting}
swift run qsb validate-dp-json manual/examples/dec_knapsack.json
swift run qsb solve-dp-json manual/examples/dec_knapsack.json
\end{lstlisting}
In macOS, Bounded Knapsack lets you edit capacity, availability, consumption,
and returns, and copy or paste the table. Summary provides metrics, bars, and
a selection table; the trace helps explain the policy.

\subsection{Stage-by-stage routing (Stagecoach)}
The \texttt{nodes} list declares a topological order: the first node is the
origin and the last the destination. Every arc must advance in that order and
have a finite, nonnegative cost. The method applies
\[
 V(i)=\min_{(i,j)}\{c_{ij}+V(j)\},\qquad V(\text{destination})=0.
\]
\path{dec_route.json} contains Origen--Norte (4), Origen--Sur (2),
Norte--Destino (3), and Sur--Destino (6). The identifiers mean Origin, North,
South, and Destination, respectively. The result is Origen--Norte--Destino,
with cost 7; choosing the cheapest initial arc would yield 8.
\begin{lstlisting}
swift run qsb validate-dp-json manual/examples/dec_route.json
swift run qsb solve-dp-json manual/examples/dec_route.json
\end{lstlisting}
A structurally valid graph may have no route to the destination; this is
detected during solution. This model does not accept cycles or backward arcs:
for other graphs, consult the networks chapter.

\subsection{Production and inventory}
The state at the beginning of a period is inventory $I_{t-1}$; the action is
integer production $P_t$. Deterministic demand $d_t$ is met with
$I_t=I_{t-1}+P_t-d_t\geq0$. Capacities bound $P_t$ and $I_t$. QSB assumes
zero initial and final inventory, with no backorders. It minimizes
\[
 \sum_t\left[K_t\mathbf{1}_{P_t>0}+c_tP_t+h_tI_t\right].
\]
Holding cost applies to the period's ending inventory, rather than its initial
inventory. The recurrence evaluates feasible production quantities and adds
the optimal future cost.

\path{dec_production.json} has demands of 2 and 3, production capacity 5,
storage capacity 3, setup cost 10, unit cost 2, and holding cost 1 in both
periods. QSB produces 5 initially, carries 3, and produces nothing in the
second period. The cost is $10+2(5)+1(3)=23$, compared with 30 if production
is set up and carried out in each period. The decisions detail initial
inventory, production, demand, ending inventory, and local cost.
\begin{lstlisting}
swift run qsb validate-dp-json manual/examples/dec_production.json
swift run qsb solve-dp-json manual/examples/dec_production.json
\end{lstlisting}
Field validation does not prove the existence of a feasible plan. Insufficient
capacities may cause solution failure. This model does not represent multiple
products, configurable initial inventory, shortages, or random demand.

\section{Use, metadata, and shared limits}
All four decision variants have native macOS editors. The tree editor allows
nodes and links to be edited; removing leaves or adding terminals maintains a
structure that is easier to review. Stagecoach and production/inventory can be
edited through the structured form or JSON; their results have native route or
per-period series views. Use Validation before Run and consult JSON when you
need to preserve every field.

The native decision backend declares \texttt{exact}: it computes the specified
criterion using numerical arithmetic within the model and its educational scale.
The three dynamic variants declare \texttt{fixtureScale}; their policies are
those obtained by the implemented recurrences, with discretization and memory
consumption growing with capacity, inventory, and the number of stages. There
is no promise of industrial-scale performance. \texttt{--backend validate}
returns validation rather than a solution; the external backend is unavailable
for these families. Save the model alongside the dynamic programming output:
its solution document preserves the title and assumptions, but not the complete
model.

\section{Exercises and checks}
\begin{enumerate}
 \item Change $P(Alta)$ from $0.4$ to $0.3$, adjusting Baja as well. First
 calculate the best decision without information, then solve again. Why can
 EVSI change even when the quality of the study remains unchanged?
 \item Calculate the maximum price of the example's study and compare it
 with EVPI. Explain which assumptions would prevent using that price in a
 real decision.
 \item In the tree, increase the cost of Invertir to 9. This branch's value
 becomes 11, and the root decision changes to Conservar, with value 12.
 \item In the knapsack, change capacity to 6. Verify a return of 10 with two
 units of A, and explain why the return/consumption ratio is insufficient to
 ensure optimality.
 \item In the route, make Sur--Destino cost 4. The route through Sur will
 cost 6. Then add a backward arc and examine the validation diagnostic.
 \item Reduce the production setup cost to 1 in both periods. Producing 2
 and 3 separately costs 12, compared with 14 when inventory is carried.
 Compare each inventory balance and the trace.
\end{enumerate}
